% Immutable cited-input supplement. Target proofs remain embargoed.
% Target: night:P4R.13; tag 02Y6.
% Frozen source SHA256 74497bfe327fae25c9cdcb17334b3bf42a9793a9df7c0a4404d97568f2f760f0.
% CompendiumRef is resolved against the cited S6_09 source edition, not silently identified with the current series statement.

% Exact cited input: series:II / night:P4.30
\begin{proposition}[P4.30]\hypertarget{tag-02X3}{}\label{night:P4.30}\label{stmt:P4.30}


Let $C$ be a compact irreducible reduced curve in the smooth
threefold $Y$ of Part~4.2. Suppose its normalization is an elliptic
curve $E$, and every point of $C$ is unibranch. Write $P$ for its
ideal in $Y$ and $P^{(2)}$ for the ideal of functions vanishing to
order at least two at the generic point of $C$. Then
\[
 \chi(\mathcal O_Y/P^{(2)})\le0.
\]
Equality forces $C$ to be smooth. The assertion uses
$c_1(TY)\cdot C=0$, which follows from $b_2(Y)=0$.
\end{proposition}

% Exact cited input: series:II / night:P4.31
\begin{proposition}[P4.31]\hypertarget{tag-02X4}{}\label{night:P4.31}\label{stmt:P4.31}


Let $S$ be an irreducible singular support in Part~4.2 whose
normalization has ruled resolution, and suppose $e(S)=0$.
In the double-horizontal-divisor case, when the generic
transverse type is the ordinary triple point $xy(x-y)=0$,
the normalization is a smooth minimal elliptic ruled surface.
Its conductor image is a smooth elliptic curve, and its
reduced conductor is an unramified cover of that curve of
degree three. Every germ along the image is analytically
the product of an ordinary plane triple point with a disc.
\end{proposition}

% Exact cited input: series:II / night:P4.33
\begin{proposition}[P4.33]\hypertarget{tag-02X6}{}\label{night:P4.33}\label{stmt:P4.33}


Let $S$ be an irreducible singular support in Part~4.2, and let
$\rho:T\to N\to S$ be its minimal normalization resolution.
Then $H^2(N,\mathcal O_N)=0$. If
$p=\dim H^0(N,R^1\rho_{N*}\mathcal O_T)$ is the sum of the
geometric genera of the normalization singularities, then
\[
 p=h^1(T,\mathcal O_T)-h^1(N,\mathcal O_N).
\]
In particular a rational resolution forces every normalization
singularity to be rational and gives $\chi(\mathcal O_N)=1$.
For the scheme-theoretic image conductor $C_{\rm sch}$ we have
\[
 \chi(\mathcal O_{C_{\rm sch}})
   =-\chi(\mathcal O_N)=h^1(N,\mathcal O_N)-1.
\]

In the ruled double-horizontal case, with section $H$ of genus
$g\ge1$, a contracted $H$ gives $h^1(N,\mathcal O_N)=0$ and
$p=g$. A noncontracted $H$ gives $p=0$ and
$h^1(N,\mathcal O_N)=g$.
\end{proposition}

% Exact cited input: series:II / night:P4R.10
\begin{proposition}[P4R.10]\hypertarget{tag-02Y3}{}\label{night:P4R.10}\label{stmt:P4R.10}


Let $S$ be an irreducible reduced fibre support in Part~4.2.
Suppose that its reduced singular curve $C$ is irreducible and
unibranch at every point, that its normalization $j:E\to C\hookrightarrow Y$
is elliptic, and that its generic transverse germ is a tacnode
$x^2-y^4=0$. Let $C_{\mathrm{sch}}$ be the image conductor scheme.
Then
\[
 \chi(\mathcal O_{C_{\mathrm{sch}}})\leq0.
\]
Equality makes $C$ smooth and makes the transverse Hessian of
the reduced defining equation have rank one at every point of $C$.
\end{proposition}

% Exact cited input: series:II / night:P4R.6
\begin{proposition}[P4R.6]\hypertarget{tag-02YI}{}\label{night:P4R.6}\label{stmt:P4R.6}


Let $S$ be an irreducible singular support in
Part~4.2, let $\nu:N\to S$ be its normalization,
and let $\rho:T\to N\to S$ use the minimal
resolution of $N$. There is an effective integral
divisor $D$ on $T$ such that
\[
 \omega_T(D)\simeq\rho^*\omega_S,
 \qquad \operatorname{Supp}D
       =\rho^{-1}(\operatorname{Sing}S)_{\rm red}.
\]
In particular $K_T+D$ is torsion, by P4.12.
The strict-transform conductor has its usual
positive integral coefficients. Every connected
exceptional configuration over a singular point
of $N$ has strictly positive coefficients in $D$.
Every such configuration meets the strict-transform
conductor.
\end{proposition}

% Exact cited input: series:II / night:P4R.7
\begin{proposition}[P4R.7]\hypertarget{tag-02YJ}{}\label{night:P4R.7}\label{stmt:P4R.7}


In Part~4.2, every irreducible singular
support whose normalization has a
resolution birationally ruled over a
curve of genus at least two has
nonnegative Euler number. The
normalization may be singular.
\end{proposition}

% Exact cited input: series:II / night:P4R.9
\begin{proposition}[P4R.9]\hypertarget{tag-02YL}{}\label{night:P4R.9}\label{stmt:P4R.9}


In Part~4.2, every irreducible singular
support whose normalization has a
resolution birationally ruled over an
elliptic curve has nonnegative Euler
number. When the horizontal part of
the pole divisor on the minimal ruled
model is reduced, the normalization
is smooth and minimal and the support
has Euler number zero.
\end{proposition}

